\subsubsection{Sidecar}
\noindent
\begin{minipage}{.49\textwidth}
    \noindent\textbf{\LARGE Sidecar\index{Sidecar}}
\end{minipage}
\hspace{.02\textwidth}
\begin{minipage}{.49\textwidth}
    \raggedleft
    \textbf{\large Brandy}\\
    \sweettartguage{0.4}\index[drinks]{Tart!40\%!Sidecar}
\end{minipage}
\vspace{.5em}\hrule\vspace{.5em}
\noindent
\textbf{\large Flavor Profile:} Tart\index[flavor-profile]{Tart}, Crisp\index[flavor-profile]{Crisp}, Orange\index[flavor-profile]{Orange}, Fruity\index[flavor-profile]{Fruity}, Hint of Grape\index[flavor-profile]{Grape}.
\vspace{.5em}\hrule\vspace{.5em}
\noindent
\textbf{\large History:} The Sidecar is a cocktail created during Prohibition but originating in France. Many popular cocktails were coming from Europe at this time as many American bartenders moved to Europe looking for work. It is difficult to know exactly who created this cocktail, but it first became popular in the 1920s. Many bartenders claimed credit for inventing the cocktail including Harry McElhone and Robert Vermier.
\vspace{.5em}\hrule\vspace{.5em}
\noindent
\begin{minipage}[t]{.49\textwidth}
    \noindent
    \textbf{\large Ingredients}
    \begin{itemize}
        \item Brandy (traditionally cognac) 60ml\index[ingredient]{Brandy}
        \item Orange Liqueur 22,5ml\index[ingredient]{Orange Liqueur}
        \item Lemon Juice 15ml\index[ingredient]{Lemon Juice}
    \end{itemize}
    \noindent
    \textbf{\large Garnish}
    \begin{itemize}
        \item Sugar Rim\index[ingredient]{Sugar}
        \item Lemon Peel\index[ingredient]{Lemon Peel}
    \end{itemize}
\end{minipage}
\hspace{.01\textwidth}\vrule\hspace{.01\textwidth}
\begin{minipage}[t]{.49\textwidth}
    \noindent
    \textbf{\large Tools/Glassware}
    \begin{itemize}
        \item Shaker
        \item Strainer \& Fine Strainer (optional)
        \item Jigger
        \item Juicer
        \item Coupe\index[glass]{Coupe} or Martini Glass\index[glass]{Martini Glass}
    \end{itemize}
\end{minipage}
\vspace{.5em}\hrule\vspace{.5em}
\noindent
\begin{minipage}{\textwidth}
    \noindent
    \textbf{\large Instructions}
    \begin{itemize}
        \item Add sugar rim to coupe or martini glass.
        \item Place all ingredients into a shaker.
        \item Add ice and shake for 8--10 seconds.
        \item Double strain into coupe or martini glass.
        \item Garnish with lemon peel.
    \end{itemize}
\end{minipage}
\vspace{.5em}\hrule\vspace{.5em}
\noindent
\textbf{\large Helpful Tip:} Orange liqueur can vary in sweetness, so add or take away to best fit your palette.
\newpage
